\clearpage\pagestyle{otMain}\setlength{\otBlocEspacement}{2.5mm}\fontsize{11}{15}\selectfont\corrigesChapitre{13}{Translations et parallélogrammes}
\begin{blocLivre}{cor-13-1}\begin{corrige}{1}{Décrire le glissement}\demarche{Reproduire le déplacement horizontal et vertical.}
Le déplacement ajoute 3 à l’abscisse et 1 à l’ordonnée. L’image est donc $\res{A'(4\,;3)}$. Le même déplacement s’applique à chaque point.
\label{fig-cor-13-1}\figureGeo[corrige]{trans-01}
\end{corrige}\end{blocLivre}
\begin{blocLivre}{cor-13-2}\begin{corrige}{2}{Construire un triangle image}\demarche{Appliquer le même déplacement aux trois sommets.}
On obtient $\res{A'(4\,;3),\ B'(6\,;3),\ C'(5\,;5)}$. Placer les trois images, puis relier dans le même ordre. Les côtés correspondants restent parallèles et de même longueur.
\label{fig-cor-13-2}\figureGeo[corrige]{trans-02}
\end{corrige}\end{blocLivre}
\begin{blocLivre}{cor-13-3}\begin{corrige}{3}{Un déplacement vers la gauche}\demarche{Relier le sens du déplacement au signe du calcul.}
Vers la gauche, on soustrait 4 à l’abscisse ; vers le bas, on soustrait 1 à l’ordonnée. Donc $\res{B'(-2\,;2)}$.
\label{fig-cor-13-3}\figureGeo[corrige]{trans-03}
\end{corrige}\end{blocLivre}
\begin{blocLivre}{cor-13-4}\begin{corrige}{4}{Conserver des mesures}\demarche{Citer la conservation adaptée à chaque mesure.}
Une translation conserve longueurs, angles et aires. Donc $\res{A'B'=5\ \mathrm{cm}}$, $\res{\widehat{B'A'C'}=42^\circ}$ et l’aire de A’B’C’ vaut $\res{9\ \mathrm{cm}^2}$.
\label{fig-cor-13-4}\figureGeo[corrige]{trans-04}
\end{corrige}\end{blocLivre}
\begin{blocLivre}{cor-13-5}\begin{corrige}{5}{Reconnaître une transformation}\demarche{Observer l’orientation du mât et du fanion.}
Le drapeau 1 a seulement glissé, sans tourner : $\res{\text{translation}}$. Le drapeau 2 est retourné par rapport à une droite verticale : $\res{\text{symétrie axiale}}$. Le drapeau 3 a tourné d’un demi-tour : $\res{\text{demi-tour}}$ ; ses segments sont parallèles à ceux de F mais de sens opposés.
\label{fig-cor-13-5}\figureGeo[corrige]{trans-05}
\end{corrige}\end{blocLivre}
\begin{blocLivre}{cor-13-6}\begin{corrige}{6}{Un parallélogramme}\demarche{Respecter l’ordre des sommets du quadrilatère.}
Les segments $[AB]$ et $[A'B']$ sont parallèles et de même longueur, car $[A'B']$ est l’image de $[AB]$ par la translation. Dans l’ordre A, B, B’, A’, ce sont deux côtés opposés. Un quadrilatère non croisé ayant deux côtés opposés parallèles et de même longueur est un $\res{\text{parallélogramme}}$.
\label{fig-cor-13-6}\figureGeo[corrige]{trans-06}
\end{corrige}\end{blocLivre}
\begin{blocLivre}{cor-13-7}\begin{corrige}{7}{Reconnaître par les diagonales}\demarche{Commencer par le milieu commun, puis ajouter le second codage.}
Figure 1 : les diagonales ont le même milieu, c’est un $\res{\text{parallélogramme}}$. Figure 2 : même milieu et diagonales perpendiculaires, c’est un $\res{\text{losange}}$. Figure 3 : même milieu et diagonales de même longueur, c’est un $\res{\text{rectangle}}$.
\label{fig-cor-13-7}\figureGeo[corrige]{trans-07}
\end{corrige}\end{blocLivre}
\begin{blocLivre}{cor-13-8}\begin{corrige}{8}{Une frise}\demarche{Compter les intervalles entre les motifs.}
Du premier au suivant : translation de 4 cm vers la droite. Du premier au cinquième, on compte quatre intervalles : $4\times4=\res{16}$ cm vers la droite. Cinq motifs ne font que quatre écarts.
\label{fig-cor-13-8}\figureGeo[corrige]{trans-08}
\end{corrige}\end{blocLivre}
\begin{blocLivre}{cor-13-9}\begin{corrige}{9}{Image d’une droite}\demarche{Utiliser les conservations.}
L’image d’une droite par une translation est une droite parallèle : $\res{(d')\parallel(d)}$. La translation conserve les angles : $(d')$ et $(\Delta')$ forment aussi un angle de $\res{35^\circ}$.
\label{fig-cor-13-9}\figureGeo[corrige]{trans-09}
\end{corrige}\end{blocLivre}
\begin{blocLivre}{cor-13-10}\begin{corrige}{10}{Deux déplacements}\demarche{Additionner séparément les déplacements sur chaque axe.}
Horizontalement : $5-3=2$ carreaux vers la droite. Verticalement : $2+4=6$ vers le haut. Le glissement total est donc de $\res{2\text{ carreaux à droite et }6\text{ vers le haut}}$.
\label{fig-cor-13-10}\figureGeo[corrige]{trans-10}
\end{corrige}\end{blocLivre}
\begin{blocLivre}{cor-13-11}\begin{corrige}{11}{Construire sans quadrillage}\demarche{Construire des parallélogrammes ayant le côté [EK].}
L’image de E est K. On construit F’ tel que EFF’K soit un parallélogramme : on trace la parallèle à $(EK)$ passant par F et l’on reporte au compas la longueur EK, dans le même sens. On procède de même pour G’. L’image cherchée est le triangle $\res{KF'G'}$, de mêmes dimensions que EFG.
\label{fig-cor-13-11}\figureGeo[corrige]{trans-11}
\end{corrige}\end{blocLivre}
\begin{blocLivre}{cor-13-12}\begin{corrige}{12}{Translation et parallélogramme}\demarche{Utiliser l’unicité du quatrième sommet.}
Notons D’ l’image de D. Par la propriété du cours, ABD’D est un parallélogramme : D’ est le quatrième sommet du parallélogramme construit sur A, B et D, dans cet ordre. Or ce quatrième sommet est C, puisque ABCD est un parallélogramme. Donc $\res{D'=C}$.
\label{fig-cor-13-12}\figureGeo[corrige]{trans-12}
\end{corrige}\end{blocLivre}
\begin{blocLivre}{cor-13-13}\begin{corrige}{13}{Un pavage par translations}\demarche{Compter les copies et convertir le déplacement en longueurs.}
Il y a $6\times4=\res{24}$ motifs, chacun d’aire 4 cm$^2$. L’aire totale vaut $\res{96\ \mathrm{cm}^2}$. Le glissement demandé est de $\res4$ cm vers la droite et $\res2$ cm vers le haut. Les translations conservent la taille du motif.
\label{fig-cor-13-13}\figureGeo[corrige]{trans-13}
\end{corrige}\end{blocLivre}
\begin{blocLivre}{cor-13-14}\begin{corrige}{14}{Deux parterres}\demarche{Utiliser la conservation des longueurs et des aires.}
DEF est l’image de ABC : mêmes longueurs, même angle droit. Chaque périmètre vaut $6+8+10=24$ m : il faut $\res{48\ \mathrm m}$ de bordure. Chaque aire vaut $\frac{6\times8}2=24$ m$^2$ : en tout $\res{48\ \mathrm m^2}$, les parterres étant disjoints. Enfin ABED est un $\res{\text{parallélogramme}}$ : A a pour image D et B pour image E.
\label{fig-cor-13-14}\figureGeo[corrige]{trans-14}
\end{corrige}\end{blocLivre}
\begin{blocLivre}{cor-13-15}\begin{corrige}{15}{Un angle transporté}\demarche{Utiliser des angles correspondants puis un angle plat.}
A a pour image E et B pour image C : ABCE est un parallélogramme. Ainsi $(CE)\parallel(BA)$. Sur le prolongement de $[BC]$ au-delà de C, l’angle formé avec $[CE)$ est le correspondant de $\widehat{ABC}$ : il mesure $70^\circ$. Comme il est adjacent à $\widehat{BCE}$ sur un angle plat, $\widehat{BCE}=180-70=\res{110^\circ}$.
\label{fig-cor-13-15}\figureGeo[corrige]{trans-15}
\end{corrige}\end{blocLivre}
